% Source: https://artofproblemsolving.com/community/c3h1398413_mathcounts_state_2011_sprint_problems
% Stable identity: aops-topic-1398413-q23
Let each piece have height $h$. The top piece is a small cone; the top two, three and four pieces together are cones similar to it with heights $2h$, $3h$, $4h$. Volumes of similar solids scale as the cube of the ratio of similarity, so if the top cone has volume $V$, the cones of heights $h,2h,3h,4h$ have volumes $V,8V,27V,64V$. The four pieces therefore have volumes
\[
V,\quad 8V-V=7V,\quad 27V-8V=19V,\quad 64V-27V=37V
\]
(as marked in the figure).

\begin{asy}
usepackage("amsmath");
size(3.8cm);
defaultpen(fontsize(10pt));
real H=4, Rb=2, k=0.22;
pair apex=(0,H);
for(int i=1;i<=4;++i){
  real y=H-i, r=Rb*i/4;
  draw(shift(0,y)*xscale(r)*yscale(k*r)*arc((0,0),1,180,360),linewidth(0.9));
  draw(shift(0,y)*xscale(r)*yscale(k*r)*arc((0,0),1,0,180),(i==4?dashed:dashed)+gray);
}
draw((-Rb,0)--apex--(Rb,0),linewidth(1));
for(int i=0;i<4;++i){ draw((Rb+0.35,H-i)--(Rb+0.55,H-i)); }
draw((Rb+0.35,0)--(Rb+0.55,0));
draw((Rb+0.45,0)--(Rb+0.45,H));
for(int i=0;i<4;++i) label("$h$",(Rb+0.45,H-i-0.5),E);
string[] vl={"$V$","$7V$","$19V$","$37V$"};
for(int i=0;i<4;++i){ real y=H-i-0.5; label(vl[i],(-Rb*(H-y)/H-0.12,y),W); }
label("\mbox{}",(min(currentpicture,user=true)+max(currentpicture,user=true))/2,black); // keeps the text after the figure black when LaTeX typesets the labels inline
\end{asy}

The largest piece is the bottom one ($37V$) and the second largest is the one above it ($19V$), so the ratio is $\boxed{\dfrac{19}{37}}$.
